% Part: second-order-logic
% Chapter: sol-and-sets
% Section: comparing-sets

\documentclass[../../../include/open-logic-section]{subfiles}

\begin{document}

\olfileid{sol}{set}{cmp}

\olsection{সেটের তুলনা}

\begin{prop}
!!{formula}~$\lforall[x][(X(x) \lif Y(x))]$ উপসেট-সম্পর্ককে
সংজ্ঞায়িত করে; অর্থাৎ, $\Sat{M}{\lforall[x][(X(x) \lif Y(x))]}[s]$
যদি এবং কেবল যদি $s(X) \subseteq s(Y)$।
\end{prop}

\begin{prop}
!!{formula}~$\lforall[x][(X(x) \liff Y(x))]$ সেটগুলির উপর অভেদ
সম্পর্ককে সংজ্ঞায়িত করে; অর্থাৎ,
$\Sat{M}{\lforall[x][(X(x) \liff Y(x))]}[s]$ যদি এবং কেবল যদি $s(X) =
s(Y)$।
\end{prop}

\begin{prop}
!!{formula}~$\lexists[x][X(x)]$ অশূন্য হওয়ার ধর্মকে সংজ্ঞায়িত করে;
অর্থাৎ, $\Sat{M}{\lexists[x][X(x)]}[s]$ যদি এবং কেবল যদি $s(X) \neq
\emptyset$।
\end{prop}

কোনো সেট~$X$, কোনো সেট~$Y$-এর চেয়ে আকারে বড় নয়, $\cardle{X}{Y}$,
যদি এবং কেবল যদি একটি !!a{injective} অপেক্ষক $f\colon X \to Y$ থাকে।
কোনো অপেক্ষক যে একৈক, এবং~$X$-এর আর্গুমেন্টগুলিতে তার মানগুলি যে~$Y$-তে
আছে, তা আমরা প্রকাশ করতে পারি। তাই সংজ্ঞাক্ষেত্রের উপসেটগুলির উপর আকারে
বড় না-হওয়ার সম্পর্কটিও সংজ্ঞায়িত করতে পারি।

\begin{prop}
নিচের সূত্রটি
\[
\lexists[u][(\lforall[x][(X(x) \lif Y(u(x)))] \land
  \lforall[x][\lforall[y][((X(x) \land X(y) \land
      \eq[u(x)][u(y)]) \lif \eq[x][y])]])]
\]
আকারে বড় না-হওয়ার সম্পর্ককে সংজ্ঞায়িত করে।
% BN-SRC-259 TARGET-CORRECTION-BEGIN
\par\smallskip\noindent\textnormal{\small\textbf{উৎস-সংশোধন (BN-SRC-259):} হিমায়িত সূত্রে~$u$-কে গোটা সংজ্ঞাক্ষেত্রে একৈক হতে বলা হয়েছে, যদিও সংজ্ঞেয় সম্পর্কটির জন্য কেবল~$X$-এ তার সীমাবদ্ধতাই একৈক হওয়া দরকার। ওই শক্ততর শর্ত সমসংখ্যক উপসেটের ক্ষেত্রেও সম্প্রসারণযোগ্য নাও হতে পারে। বাংলা সূত্রে একৈকতার পূর্বগে~$X(x) \land X(y)$ যোগ করে যথাযথ সীমাবদ্ধতা পুনরুদ্ধার করা হয়েছে।} % BN-SRC-259 TARGET-CORRECTION-END
\end{prop}

দুটি সেটের আকার একই, অথবা তারা ``সমসংখ্যক,'' $\cardeq{X}{Y}$, যদি এবং
কেবল যদি একটি !!a{bijective} অপেক্ষক~$f\colon X \to Y$ থাকে।

\begin{prop}
নিচের সূত্রটি
\begin{multline*}
  \lexists[u][(\lforall[x][(X(x) \lif Y(u(x)))] \land {}\\
    \lforall[x][\lforall[y][((X(x) \land X(y) \land
      \eq[u(x)][u(y)]) \lif \eq[x][y])]]
  \land {} \\\lforall[y][(Y(y) \lif \lexists[x][(X(x)
      \land \eq[y][u(x)])])])]
\end{multline*}
সমসংখ্যক হওয়ার সম্পর্ককে সংজ্ঞায়িত করে।
% BN-SRC-260 TARGET-CORRECTION-BEGIN
\par\smallskip\noindent\textnormal{\small\textbf{উৎস-সংশোধন (BN-SRC-260):} হিমায়িত সূত্রটি~$X$ থেকে~$Y$-তে দ্বৈকতার বদলে~$u$-কে গোটা সংজ্ঞাক্ষেত্রে একৈক হতে বলে। যেমন, অসীম সংজ্ঞাক্ষেত্রে~$X$ একটি এক-উপাদান-বর্জিত উপসেট ও~$Y$ গোটা সংজ্ঞাক্ষেত্র হলে তারা সমসংখ্যক হলেও এমন বিশ্বব্যাপী একৈক সম্প্রসারণ অসম্ভব। বাংলা সূত্রে একৈকতা~$X$-এর উপাদানযুগলের মধ্যে সীমাবদ্ধ করা হয়েছে।} % BN-SRC-260 TARGET-CORRECTION-END
\end{prop}

এই !!{formula}s-কে আমরা যথাক্রমে $X \subseteq Y$, $X = Y$, $X \neq
\emptyset$, $\cardle{X}{Y}$ ও~$\cardeq{X}{Y}$ দিয়ে সংক্ষেপ করব। (এতে
সামান্য বিভ্রান্তি হতে পারে, কারণ $X$ ও~$Y$ সেট নিয়ে অনানুষ্ঠানিকভাবে
কথা বলার সময়ও আমরা একই সংকেত ব্যবহার করি---কিন্তু এখানে সংকেতগুলি
এক-স্থানীয় সম্পর্ক-চলরাশি $X$ ও~$Y$-সম্পৃক্ত দ্বিতীয়-ক্রমের
যুক্তিবিদ্যার !!{formula}s-এর সংক্ষিপ্ত রূপ।)

\begin{prop}
!!{sentence}~$\lforall[X][\lforall[Y][((\cardle{X}{Y} \land
    \cardle{Y}{X}) \lif \cardeq{X}{Y})]]$ বৈধ।
\end{prop}

\begin{proof}
!!{sentence}-টি কোনো !!a{structure}~$\Struct{M}$-এ পরিতৃপ্ত হয় যদি,
যেকোনো উপসেট $X \subseteq \Domain{M}$ ও~$Y \subseteq \Domain{M}$-এর
ক্ষেত্রে, $\cardle{X}{Y}$ এবং $\cardle{Y}{X}$ হলে
$\cardeq{X}{Y}$ হয়। কিন্তু এটি \emph{যেকোনো} সেট $X$ ও~$Y$-এর জন্যই
সত্য---এটাই শ্র্যোডার--বার্নস্টাইন উপপাদ্য।
\end{proof}


\end{document}
